\documentclass{article}

\usepackage{graphicx}
\usepackage{listings}

\begin{document}

\begin{center}


\includegraphics[width=3cm]{iiitb_comet_logo.jpg}

\vspace{0.4cm}

IIIT Bangalore

\vspace{0.3cm}

Name: Srinidhi A

ID: COMETFWC053

\vspace{0.5cm}

Question 8

\end{center}

\section*{Question}

Simplify the Boolean function using the Karnaugh Map shown.

\section*{Analysis}

A 4-variable Karnaugh map is used to simplify Boolean expressions.

Variables used:

$P, Q, R, S$

The map is arranged in Gray code order.

Rows represent $RS$ and columns represent $PQ$.

By grouping adjacent 1s in the K-map we obtain the simplified Boolean expression.

\section*{Python Code}

\begin{lstlisting}
import numpy as np
import matplotlib.pyplot as plt

kmap = np.array([
[0,1,0,0],
[0,1,1,1],
[1,1,1,0],
[0,0,1,0]
])

print("K-map Matrix:")
print(kmap)

plt.figure()
plt.imshow(kmap)

for i in range(4):
    for j in range(4):
        plt.text(j,i,kmap[i][j],ha='center',va='center')

plt.title("4 Variable Karnaugh Map")
plt.xlabel("PQ (00 01 11 10)")
plt.ylabel("RS (00 01 11 10)")

plt.colorbar()

plt.savefig("graph.png")

plt.show()
\end{lstlisting}

\section*{Graph}

\begin{center}

\includegraphics[width=10cm]{graph.png}

\end{center}

\section*{Conclusion}

The Karnaugh Map visualization helps identify groups of adjacent 1s. These groups allow simplification of the Boolean expression, reducing the number of logic gates required in the final circuit implementation.

\end{document}
